\chapter{Допуск Real-ориентированного голоморфного цикла}
\label{ch:v8-horizontal-phase-real-oriented-cycle-admission-gate}

Предыдущий гейт нашёл три обычных цикла, но показал сокращение их
горизонтального заряда. Остаётся проверить более сильное предположение:
может ли вещественная структура превратить сопряжённые стрелки в независимый
ориентированный голоморфный путь.

Пусть $N$ --- повышающая матрица полной одиннадцатирёберной опоры в порядке
вершин $(Q_L,L_L,X_L,Y_L\mid u_R,d_R,e_R,X_R,Y_R)$. Она имеет блочный вид
\begin{equation}
 N=\begin{pmatrix}0_{4}&0\\A_{5\times4}&0_{5}\end{pmatrix},
 \qquad N^2=0.
 \label{eq:v8-real-oriented-raising-nilpotence}
\end{equation}
Следовательно, чисто голоморфный двуслойный колчан не содержит даже пути
длины два, а значит, не содержит направленного цикла.

Физическая Real-структура добавляет обратный блок не как новое поле, а как
эрмитово сопряжение:
\begin{equation}
 N^{\rm rev}=N^*,
 \qquad
 D_{\mathbb R}=N+N^*.
 \label{eq:v8-real-oriented-real-completion}
\end{equation}
Поэтому двадцатимерный комплексный носитель переноса остаётся
сорокамерным вещественным носителем,
\begin{equation}
 \dim_{\mathbb R}\mathcal T_{\mathbb R}=2\cdot20=40,
 \label{eq:v8-real-oriented-physical-dimension}
\end{equation}
а объявление обратной копии независимой дало бы
\begin{equation}
 \dim_{\mathbb R}(\mathcal T\oplus\mathcal T^{\rm rev})=4\cdot20=80,
 \qquad \Delta\dim_{\mathbb R}=40.
 \label{eq:v8-real-oriented-independent-reverse-dimension}
\end{equation}
Это было бы новым полевым сектором, а не следствием Real-завершения.

Для градуировки $\Gamma=\operatorname{diag}(I_4,-I_5)$ выполнено
\begin{equation}
 \Gamma D_{\mathbb R}+D_{\mathbb R}\Gamma=0.
 \label{eq:v8-real-oriented-odd-grading}
\end{equation}
Отсюда нечётные следовые слова исчезают. Более того, горизонтальная фаза
действует матрицей
$G(z)=\operatorname{diag}(I_4,z^4,z^3I_4)$ и даёт точное подобие
\begin{equation}
 D_{\mathbb R}(z)=G(z)D_{\mathbb R}(1)G(z)^{-1}.
 \label{eq:v8-real-oriented-phase-similarity}
\end{equation}
Обратные стрелки автоматически получают противоположные веса $-4$ и $-3$.
На полном представителе опоры первые моменты равны
\begin{equation}
 \bigl(\Tr D,\Tr D^3,\Tr D^5\bigr)=(0,0,0),
 \qquad
 \bigl(\Tr D^2,\Tr D^4,\Tr D^6\bigr)=(22,110,682),
 \label{eq:v8-real-oriented-trace-moments}
\end{equation}
и не зависят от $z$. В вершинной записи сокращение имеет вид
\begin{equation}
 (4,3,3,3,3)+(-4,-3,-3,-3,-3)=0.
 \label{eq:v8-real-oriented-charge-cancellation}
\end{equation}

Наконец, нельзя исправить ситуацию фазовым переопределением формальной
обратной стрелки. Если $b=\alpha a^*$ и $b^*=\beta a$, то сохранение
ориентированного предпроективного знака требует
\begin{equation}
 \alpha\beta=-1,
 \label{eq:v8-real-oriented-formal-sign}
\end{equation}
тогда как положительная инволюция требует
\begin{equation}
 \beta=\overline\alpha,
 \qquad
 \alpha\beta=|\alpha|^2\geq0.
 \label{eq:v8-real-oriented-positive-involution}
\end{equation}
Условия несовместимы. Это воспроизводит на нынешнем полном носителе ранний
запрет производного отображения момента: формальная ориентация требует
дополнительной симплектической поляризации.

\begin{theorem}[Real-завершение без голоморфного цикла]
Чисто повышающая часть полного двуслойного колчана удовлетворяет $N^2=0$.
Real-обратные стрелки однозначно равны сопряжённым и несут противоположные
фазовые веса. Полный нечётный оператор меняется вдоль горизонтальной
окружности подобием, поэтому все его спектрально-следовые слова фазонезависимы.
\end{theorem}

\begin{nogocriterion}[No-go независимой Real-обратной стрелки]
Нельзя считать $N^*$ новым голоморфным полем: это удваивает вещественный
носитель переноса с $40$ до $80$ направлений. Нельзя также вставить минус
в формальную обратную стрелку, сохраняя положительную инволюцию. Следующий
гейт может рассматривать комплексную симплектическую поляризацию только как
явно проверяемое расширение архитектуры с собственной формой и мерой.
\end{nogocriterion}

\begin{protocol}
\begin{enumerate}
 \item чисто повышающий квадрат: \textbf{$N^2=0$};
 \item направленные голоморфные циклы: \textbf{$0$};
 \item Real-обратная стрелка независима: \textbf{нет};
 \item физическая размерность переноса: \textbf{$40$ вещественных};
 \item размерность при независимой обратной копии: \textbf{$80$};
 \item нечётные следовые моменты: \textbf{$(0,0,0)$};
 \item чётные моменты степеней $2,4,6$: \textbf{$(22,110,682)$};
 \item действие горизонтальной фазы: \textbf{подобие};
 \item суммарный заряд Real-пары: \textbf{$0$};
 \item совместимость формального минуса с положительной инволюцией:
 \textbf{нет};
 \item Real-ориентированный цикл поднимает моду: \textbf{нет};
 \item следующий гейт: \textbf{допуск комплексной симплектической
 поляризации}.
\end{enumerate}
\end{protocol}

Машинный сертификат сохранён в
\path{s2t_v8_horizontal_phase_real_oriented_cycle_admission_gate_results.json}.

\begin{thebibliography}{9}
\bibitem{v8-real-oriented-preprojective}
W.~Crawley-Boevey,
\emph{Poisson structures on moduli spaces of representations},
J. Algebra 325 (2011), 205--215; arXiv:math/0502301.
\bibitem{v8-real-oriented-krajewski}
T.~Krajewski, \emph{Classification of Finite Spectral Triples},
J. Geom. Phys. 28 (1998), 1--30; arXiv:hep-th/9701081.
\end{thebibliography}